\documentclass[11pt,a4paper]{article}
\usepackage[utf8]{inputenc}
\usepackage[spanish]{babel}
\usepackage{geometry}
\usepackage{booktabs}
\usepackage{array}
\usepackage{multirow}
\usepackage{graphicx}
\usepackage{xcolor}
\usepackage{hyperref}
\usepackage{fancyhdr}
\usepackage{pifont}
\newcommand{\cmark}{\ding{51}} % tick
\newcommand{\xmark}{\ding{55}} % cruz

\geometry{
    a4paper,
    left=2.5cm,
    right=2.5cm,
    top=3cm,
    bottom=3cm,
    headheight=14pt
}

\addto\captionsspanish{\renewcommand{\tablename}{Tabla}}

\pagestyle{fancy}
\fancyhf{}
\rhead{Laboratorio 1 - EXIST 2025}
\lhead{Detección de sexismo}
\rfoot{\thepage}

\definecolor{azuloscuro}{RGB}{0,51,102}
\definecolor{azulclaro}{RGB}{51,102,153}

\hypersetup{
    colorlinks=true,
    linkcolor=azuloscuro,
    citecolor=azuloscuro,
    urlcolor=azulclaro
}

\title{%
    \huge Clasificación de sexismo en memes\\[0.2cm]
    \large{MIARFID - ALC - Laboratorio 2}
}
\author{Pablo Segovia Martínez - Shiyi Cheng}
\date{\today}

\begin{document}

\maketitle

\section{Introducción}

Detección automática de contenido sexista en memes usando el dataset EXIST. Clasificación binaria (sexista/no sexista) mediante modelos de embeddings de texto, de vision y de lenguaje.

\section{Metodología}

\subsection{Preprocesamiento}

    Se dispone de un dataset con imágenes, su texto extraído, su descripción y datos multimodales tales como EEG, ET y HR. En primer lugar, hemos realizado un análisis previo para determinar cuáles son las variables multimodales significativas para su posible uso posterior. Las variables seleccionadas han sido las siguientes:
    \begin{itemize}\setlength{\itemsep}{0pt}\setlength{\parskip}{0pt}\setlength{\topsep}{2pt}
    	\item \texttt{ET\_reaction\_time}
    	\item \texttt{ET\_fixations\_count}
    	\item \texttt{ET\_fixations\_duration\_mean\_ns}
    	\item \texttt{ET\_fixations\_duration\_std\_ns}
    	\item \texttt{ET\_saccades\_count}
    	\item \texttt{ET\_saccades\_duration\_mean\_ns}
    	\item \texttt{ET\_blinks\_count}
    	\item \texttt{ET\_blinks\_duration\_min\_ns}
    	\item \texttt{ET\_blinks\_duration\_mean\_ns}
    	\item \texttt{HR\_garmin\_hr\_std}
    	\item \texttt{HR\_garmin\_hr\_mean}
    	\item \texttt{HR\_garmin\_hr\_max}
    	\item \texttt{HR\_garmin\_hr\_min}
    \end{itemize}
    
	Además, la descripción de la imagen junto con el texto extraído se ha utilizado como entrada para los modelos de solo texto.
	
\subsection{Modelos evaluados}
\begin{itemize}
	\item Lineales: MLP de caraterísticas sensoriales.
    \item Modelos de embeddings FT: F2LLM-4B.
    \item Modelos de visión FT: DinoV3-vitl16 y convnextv2-large-22k-224.
    \item Ensamble de probabilidades de MLP, F2LLM y DinoV3.
    \item Multimodal que concatena embeddings de DinoV3, F2LLM y del MLP sensorial para ser procesados por un MLP.
    \item Vision LLM con y sin FT: InternVL3\_5-8B, Pixtral-12b y Qwen3-VL-8B-Instruct.
\end{itemize}

\section{Resultados}
\begin{table}[h]
	\centering
	\small
	\begin{tabular}{@{}lcccc@{}}
		\toprule
		\textbf{Modelo} & \textbf{Entrada} & \textbf{FT} & \textbf{Acc} & \textbf{F1-Macro} \\
		\midrule
		\multirow{1}{*}{MLP sensorial}
		& sensorial          & \xmark       &    0.57          &     0.57         \\
		\midrule
		\multirow{2}{*}{F2LLM-4B}
		& texto          & \cmark       &   0.72          &    0.69          \\
		& descripción + texto          & \cmark       &   0.75           &     0.72        \\
		\midrule
		\multirow{1}{*}{DinoV3}
		& imagen          & \cmark       &    0.60          &     0.39         \\
		\midrule
		\multirow{1}{*}{ConvNext}
		& imagen             & \cmark       &     0.60         &    0.37        \\
		\midrule
		\multirow{1}{*}{Ensamble}
		& probabilidades             & \xmark       &     0.74         &    0.69        \\
		\midrule
		\multirow{1}{*}{Multimodal MLP}
		& embeddings             & \xmark       &     0.78         &    \textbf{0.77}        \\
		\midrule
		\multirow{2}{*}{InternVL 3.5 8B}
		& imagen + descripción + texto    &  \xmark       &     0.77         &    0.75        \\
		& imagen + descripción + texto    &  \cmark       &     0.73         &    0.70        \\
		\midrule
		\multirow{2}{*}{Pixtral 12B}
		& imagen + descripción + texto    &  \xmark       &     0.77         &    0.75        \\
		& imagen + descripción + texto    &  \cmark       &     0.60         &    0.37        \\
		\midrule
		\multirow{2}{*}{Qwen3VL 8B}
		& imagen + descripción + texto    &  \xmark       &     0.76         &    0.72        \\
		& imagen + descripción + texto    &  \cmark       &     \textbf{0.79}         &    \textbf{0.77}        \\
		\bottomrule
	\end{tabular}
	\caption{Comparativa de todos los modelos utilizados.}
\end{table}

\section{Conclusiones}
	Aquellos modelos que no son multimodales por si solos proporcionan resultados pobres, pero si se combinan se obtienen muy buenos resultados, como es el caso de \textbf{Multimodal MLP}. Por otra parte, modelos multimodales basados en Vision LLM con \textit{zero-shot} generan buenos resultados, pero si se ajustan a la tarea pueden igualar/mejorar a modelos más \textit{adhocs} como el Multimodal MLP, como es el caso de \textbf{Qwen3VL 8B}.

\end{document}
